\section{Problem Statement}

\textbf{The SART Problem.}
We study supervised training of language models on reasoning tasks. Let $\mathcal{D} = \{(x_i, y_i)\}_{i=1}^{N}$ be a training dataset where $x_i$ is an input reasoning problem and $y_i$ is the corresponding output comprising intermediate reasoning steps and a final answer. Let $\mathcal{V} = \{(x_j, y_j)\}_{j=1}^{M}$ be a small held-out validation set drawn from a distribution that rewards genuine reasoning, and let $\theta$ denote the parameters of a language model $f_\theta$. Standard training minimizes the empirical risk:
\begin{equation}
    \min_{\theta} \; \frac{1}{N} \sum_{i=1}^{N} \ell(\theta;\, x_i, y_i),
\end{equation}
where $\ell$ is the token-level negative log-likelihood loss. This objective does not distinguish between samples that promote genuine logical inference and those that achieve low loss via superficial patterns---keyword-answer correlations, answer token memorization, or template matching. We call such samples \emph{shortcut reasoning samples}.

Formally, let $\mathbf{g}_s = \nabla_\theta\, \ell(\theta; s)$ denote the gradient of a training sample $s = (x, y)$, and let
\begin{equation}
    \mathbf{g}_{\mathcal{V}} = \nabla_\theta \; \frac{1}{|\mathcal{V}|} \sum_{v \in \mathcal{V}} \ell(\theta;\, v)
\end{equation}
be the validation gradient, reflecting parameter updates that improve generalizable reasoning. A sample $s$ is a shortcut reasoning sample if $\mathbf{g}_s$ (i) exhibits low alignment with $\mathbf{g}_{\mathcal{V}}$, indicating the update does not transfer to genuine reasoning, and (ii) concentrates gradient norm on final answer tokens rather than intermediate reasoning tokens, indicating the model learns \emph{what to answer} rather than \emph{how to reason}. Critically, such samples are often correctly labeled and efficiently reduce training loss, rendering them undetectable by conventional data quality, label-noise, or scalar-loss metrics.

\textbf{Shortcut-Aware Reasoning Training (SART)} formulates the problem as a two-step gradient-geometry detection-then-correction task. Concretely, we pursue two objectives: 1) \textbf{detection}: rank training samples by gradient-derived shortcut tendency, given by a \textsc{ShortcutScore} $S(s)$ that combines non-transfer alignment with answer-token concentration; 2) \textbf{conditional correction}: for samples flagged as shortcut-prone, project out only the component of $\mathbf{g}_s$ that genuinely opposes $\mathbf{g}_{\mathcal{V}}$, leaving validation-aligned directions intact. Formally, with $\mathcal{S}^\star = \{s : S(s) > \tau_S^\star\}$ the flagged set and $\mathbf{g}_s'$ the corrected per-sample gradient, SART solves
\begin{equation}
    \min_{\theta} \; \frac{1}{N}\sum_{i=1}^{N} \ell(\theta;\, x_i, y_i)
    \quad \text{subject to} \quad
    \mathbf{g}_s'^{\top} \mathbf{g}_{\mathcal{V}} \;\geq\; 0 \;\;\text{for all } s \in \mathcal{S}^\star,
\end{equation}
where the constraint is realized by a conflict-only PCGrad-style projection on the per-sample gradient before batch aggregation; samples with $s \notin \mathcal{S}^\star$ contribute their original gradient at unit weight. The framework operates purely at the level of training dynamics, requiring neither architectural modifications nor shortcut-group / multi-environment annotations.
